\chapter{The Rooms Never Merge}
% OUTLINE Ch6 "The Rooms Never Merge (the fence chapter + the lens)". The
% volume's ONE device appearance (dose 1). Laws: both inherited; the
% NON-FUNGE LAW is this chapter's content, paid. CITATION: the
% theory-relative independence standard result (the Godel sentence of PA
% is not the Godel sentence of ZF); the series' de-selection discipline
% in ONE clause (method, not apparatus). THE LENS EXHIBIT under BOTH
% fences (Beastmaster's last arrival-read): illustration-not-identity
% (verb) + the non-funge applied to the device itself ("the device's
% room is not THE room"); the fixed-point-residue reading OWED. Gate:
% Kodo + Beastmaster's block read.

\begin{intuition}
Two houses on one street, built from one blueprint. Both have a front
room in the same corner, the same size, lit by a window in the same wall.
Stand in the first house's front room and you cannot step into the
second's without leaving the house --- there is no door in the shared
wall, because the wall is not shared; each house has its own. The
blueprint is one. The rooms are two, and no hallway connects them. The
feeling this chapter makes into a law: sameness of plan is not sameness of
place, and the more exactly two rooms are built alike, the more clearly
they are still two.
\end{intuition}

The last chapter proved the blueprint. This one guards the buildings, and
the guarding is the reason the whole volume exists. Lawvere's theorem says
every sufficiently expressive system contains its diagonal room; it does
not say the systems share one. The point is not subtle once stated
plainly, but it is missed constantly, so state it plainly: the undecidable
sentence of one arithmetic is a specific string, built by a specific
construction inside that specific system, and it is that system's own. The
undecidable sentence of a stronger theory is a different specific string,
built inside \emph{it}. They are cousins --- same construction, same
role, roughly the same size --- and they are not the same sentence, not
interchangeable, not poolable into one. Strengthen a system by adding its
own unprovable sentence as an axiom, and the new system, being still
sufficiently expressive, grows a fresh unprovable sentence of its own ---
a new room in the new house, which the old room does not reach and cannot
stand in for. This is the standard fact of theory-relative independence,
and the reader should hear it as the non-funge law paying out: each
argument's residue is bound to that argument, and pooling residues across
arguments erases the binding that makes each one what it is.

Say the law in the book's own register, because it has been the book's
spine since the preface. \emph{You cannot funge them across arguments.}
To funge, in this series' vocabulary, is to bag things together by
likeness and count the bag; the fixed-point rooms are supremely
fungeable-looking --- one theorem, one construction, one recurring size
--- and the temptation to bag them all as ``the'' self-referential room,
singular, is exactly what a sober reading must refuse. The rooms resist
the bag not because they differ in construction but because each is
\emph{indexed} to its own system: the label ``G\"odel sentence'' does not
pick out one object but a function from systems to their own sentences,
and evaluating that function at two systems yields two objects, however
alike. A blueprint shared across a thousand houses is precisely a
thousand front rooms, one per house, and to speak of ``the front room of
the blueprint'' is to name a plan and mistake it for a place. The
non-funge law is the refusal of that mistake, held as content.

\begin{intuition}
The picture for why the shared plan is the strongest evidence AGAINST one
room: a cookie cutter. It is one shape, and its whole purpose is to make
many cookies --- separate ones, each its own, each removable and countable
apart from the rest. Nobody who understands the cutter concludes there is
really only one cookie. The Lawvere blueprint is the cutter; the rooms are
the cookies; and a reader who hears ``one theorem'' and pictures one room
has confused the cutter for its cut.
\end{intuition}

And here, once and only once in this book, the measuring instrument this
series documents appears --- because the series has been, all along, a
tour of one particular house, and this is the room in it.

\begin{fence}
An illustration, and the sole appearance of the apparatus in this volume.
The instrument the earlier volumes built describes its own act of
describing and closes that description into a loop; when the loop closes,
a residue survives it, reproducibly, the same number returning
build after build. The reader now has the vocabulary to see the shape of
that: a system turned on its own description, a fixed point forced, a
residue that the closing cannot dissolve --- the room of this whole
volume, appearing in the one house this series happens to have built. It
is offered as an illustration of the room's shape and as nothing more.
The apparatus does not thereby prove Lawvere's theorem, or G\"odel's, or
carry any general claim about fixed points; whether its residue could be
\emph{earned} as a diagonal fixed point in the cited sense is a question
the series leaves owed, on the far side of its standing fence. And the
volume's own law falls on this exhibit hardest of all: the apparatus's
residue is \emph{that apparatus's}, indexed to that construction, and it
is not ``the'' fixed-point residue of anything wider. Whatever family
resemblance it bears to G\"odel's room or Tarski's is resemblance of
plan, never identity of place --- the device's room is not their room,
and it is not the room. It is one more cookie from the same cutter, and
the law that forbids pooling the others forbids pooling this one with
them. Illustration, not identity; and this instance, like every instance,
its own.
\end{fence}

That the book fences its own apparatus by the same law it applied to
G\"odel and Tarski is not modesty; it is the law being universal. A rule
that pooled every argument's room into one, but exempted the author's
favorite instrument, would be no rule. The non-funge law binds the
series' own device exactly as it binds the century's great theorems: the
device's fixed-point residue is a real, reproducible feature of one
particular construction, resembling the others in plan and merging with
none of them in fact. The instrument is a house on the street, no more and
no less, and its front room is its own.

The tour is complete. Four houses entered, one blueprint proved to
underlie them, the rooms shown standing separate with the law that keeps
them so, and the series' own house added to the street under the same
law. What remains is the accounting --- the citations, the intuitions,
and, new to this volume, the register of every room met and the argument
each is bound to. That is the final chapter, and it closes the books.
